\documentclass[10pt,a4paper]{amsart}
\usepackage{amsmath,amssymb,amsthm,graphicx}
\usepackage{verbatim}
\newenvironment{reference}{\comment}{\endcomment}
%\newenvironment{reference}{}{}
\newenvironment{slogan}{\comment}{\endcomment}
\newenvironment{history}{\comment}{\endcomment}
% For commutative diagrams we use Xy-pic
\usepackage[all]{xy}

% We use 2cell for 2-commutative diagrams.
\xyoption{2cell}
% Theorem environments.
%
\theoremstyle{plain}
\newtheorem{theorem}[subsection]{Theorem}
\newtheorem{proposition}[subsection]{Proposition}
\newtheorem{lemma}[subsection]{Lemma}

\theoremstyle{definition}
\newtheorem{definition}[subsection]{Definition}
\newtheorem{example}[subsection]{Example}
\newtheorem{exercise}[subsection]{Exercise}
\newtheorem{situation}[subsection]{Situation}

\theoremstyle{remark}
\newtheorem{remark}[subsection]{Remark}
\newtheorem{remarks}[subsection]{Remarks}

\numberwithin{equation}{subsection}

% Macros
%
\def\lim{\mathop{\mathrm{lim}}\nolimits}
\def\colim{\mathop{\mathrm{colim}}\nolimits}
\def\Spec{\mathop{\mathrm{Spec}}}
\def\Hom{\mathop{\mathrm{Hom}}\nolimits}
\def\Ext{\mathop{\mathrm{Ext}}\nolimits}
\def\SheafHom{\mathop{\mathcal{H}\!\mathit{om}}\nolimits}
\def\SheafExt{\mathop{\mathcal{E}\!\mathit{xt}}\nolimits}
\def\Sch{\mathit{Sch}}
\def\Mor{\mathop{\mathrm{Mor}}\nolimits}
\def\Ob{\mathop{\mathrm{Ob}}\nolimits}
\def\Sh{\mathop{\mathit{Sh}}\nolimits}
\def\NL{\mathop{N\!L}\nolimits}
\def\CH{\mathop{\mathrm{CH}}\nolimits}
\def\proetale{{pro\text{-}\acute{e}tale}}
\def\etale{{\acute{e}tale}}
\def\QCoh{\mathit{QCoh}}
\def\Ker{\mathop{\mathrm{Ker}}}
\def\Im{\mathop{\mathrm{Im}}}
\def\Coker{\mathop{\mathrm{Coker}}}
\def\Coim{\mathop{\mathrm{Coim}}}

% Boxtimes
%
\DeclareMathSymbol{\boxtimes}{\mathbin}{AMSa}{"02}

%
% Macros for moduli stacks/spaces
%
\def\QCohstack{\mathcal{QC}\!\mathit{oh}}
\def\Cohstack{\mathcal{C}\!\mathit{oh}}
\def\Spacesstack{\mathcal{S}\!\mathit{paces}}
\def\Quotfunctor{\mathrm{Quot}}
\def\Hilbfunctor{\mathrm{Hilb}}
\def\Curvesstack{\mathcal{C}\!\mathit{urves}}
\def\Polarizedstack{\mathcal{P}\!\mathit{olarized}}
\def\Complexesstack{\mathcal{C}\!\mathit{omplexes}}
% \Pic is the operator that assigns to X its picard group, usage \Pic(X)
% \Picardstack_{X/B} denotes the Picard stack of X over B
% \Picardfunctor_{X/B} denotes the Picard functor of X over B
\def\Pic{\mathop{\mathrm{Pic}}\nolimits}
\def\Picardstack{\mathcal{P}\!\mathit{ic}}
\def\Picardfunctor{\mathrm{Pic}}
\def\Deformationcategory{\mathcal{D}\!\mathit{ef}}
\usepackage{cleveref}
\crefname{originallemma}{original lemma}{original lemmas}
\Crefname{originallemma}{Original lemma}{Original lemmas}
\makeatletter
\newlabel{lemma-composition-finite-type}{{lemma-composition-finite-type}{source}{}{sourcepointer-1}{}}
\newlabel{lemma-composition-finite-type@cref}{{[originallemma][0][]lemma-composition-finite-type}{[1][0][]source}{}{}{}}
\makeatother
\crefname{originalreference}{original reference}{original references}
\Crefname{originalreference}{Original reference}{Original references}
\makeatletter
\newlabel{morphisms-lemma-identify-finite-type-points}{{morphisms-lemma-identify-finite-type-points}{source}{}{sourcepointer-2}{}}
\newlabel{morphisms-lemma-identify-finite-type-points@cref}{{[originalreference][0][]morphisms-lemma-identify-finite-type-points}{[1][0][]source}{}{}{}}
\makeatother
\crefname{originalreference}{original reference}{original references}
\Crefname{originalreference}{Original reference}{Original references}
\makeatletter
\newlabel{spaces-morphisms-lemma-etale-locally-finite-presentation}{{spaces-morphisms-lemma-etale-locally-finite-presentation}{source}{}{sourcepointer-3}{}}
\newlabel{spaces-morphisms-lemma-etale-locally-finite-presentation@cref}{{[originalreference][0][]spaces-morphisms-lemma-etale-locally-finite-presentation}{[1][0][]source}{}{}{}}
\makeatother
\crefname{originalreference}{original reference}{original references}
\Crefname{originalreference}{Original reference}{Original references}
\makeatletter
\newlabel{spaces-morphisms-lemma-finite-presentation-finite-type}{{spaces-morphisms-lemma-finite-presentation-finite-type}{source}{}{sourcepointer-4}{}}
\newlabel{spaces-morphisms-lemma-finite-presentation-finite-type@cref}{{[originalreference][0][]spaces-morphisms-lemma-finite-presentation-finite-type}{[1][0][]source}{}{}{}}
\makeatother
\crefname{originalreference}{original reference}{original references}
\Crefname{originalreference}{Original reference}{Original references}
\makeatletter
\newlabel{spaces-morphisms-lemma-finite-type-points-morphism}{{spaces-morphisms-lemma-finite-type-points-morphism}{source}{}{sourcepointer-5}{}}
\newlabel{spaces-morphisms-lemma-finite-type-points-morphism@cref}{{[originalreference][0][]spaces-morphisms-lemma-finite-type-points-morphism}{[1][0][]source}{}{}{}}
\makeatother
\crefname{originalreference}{original reference}{original references}
\Crefname{originalreference}{Original reference}{Original references}
\makeatletter
\newlabel{spaces-morphisms-lemma-identify-finite-type-points}{{spaces-morphisms-lemma-identify-finite-type-points}{source}{}{sourcepointer-6}{}}
\newlabel{spaces-morphisms-lemma-identify-finite-type-points@cref}{{[originalreference][0][]spaces-morphisms-lemma-identify-finite-type-points}{[1][0][]source}{}{}{}}
\makeatother
\title{Two equivalent descriptions of a point of finite type of an algebraic stack}
\author{The Stacks Project Authors}
\date{stacks-project:stacks-morphisms.tex:local-snapshot-20261012; GNU FDL 1.2 or later}
\begin{document}
\maketitle
\noindent\textit{Editorial note.} All delivered mathematical text is the author's own, copied from the pinned Stacks Project chapter source sources/stacks-algebraic-stacks-morphisms/pinned-original-source.tex; no statement or proof step was rewritten, reordered or removed.\par\medskip
\noindent\textit{Editorial note.} Preamble provenance: the Stacks Project keeps one single shared preamble for every chapter; this card takes it directly from the staged sources/stacks-algebraic-stacks-morphisms/preamble.tex as three verbatim spans (lines 8-12, 14-18 and 154-217), so no derived file was written and nothing was added or redefined. The wrapper class and the base packages are supplied by the assembler.\par\medskip
\noindent\textit{Editorial note.} Bibliography provenance: neither the counted statement nor its proof contains a citation (this lane rescans the whole span before writing the card), so the card carries no bibliography and no entry was reconstructed.\par\medskip
\noindent\textit{Editorial note.} Selection provenance: this card was opened under the file-based control contract (CONTROL.md), with the mechanical pre-card gate run on its statement and proof interval before the card was written.\par\medskip
\noindent\textit{Allowed standard prerequisites (editorial).}\begin{itemize}
\item The statement that a smooth morphism of algebraic spaces is locally of finite presentation and hence locally of finite type. Invoked at the start of the counted proof (source lines 3668-3672) as a cross-chapter pointer. Makes the presenting morphism locally of finite type.
\item The composition permanence of locally of finite type morphisms. Invoked inside the counted proof (source lines 3673-3674). Turns the residue field point through the presentation into a locally finite type morphism.
\item The existence of finite type points on a nonempty algebraic space of locally finite type over a field. Invoked inside the counted proof (source lines 3678-3681) as a cross-chapter pointer. Produces a finite type point of the fibre of the presentation.
\item The statement that the image of a finite type point under a morphism of finite type is again a finite type point, and the description of finite type points of a scheme through closed points. Invoked near the end of the counted proof (source lines 3682-3690) as cross-chapter pointers. Shows that the point can be assumed to be a closed point of the presentation.
\end{itemize}\medskip
\ifdefined\hypertarget\hypertarget{sourcepointer-1}{}\fi
\noindent\textit{Source reference.} Original lemma lemma-composition-finite-type: sources/stacks-algebraic-stacks-morphisms/pinned-original-source.tex, source lines 3488--3488. Navigation pointer only: the counted text refers to the Stacks item lemma-composition-finite-type, which is an existing labeled result of this same chapter and is not part of this card. The pointer carries the target's own label as its number and its definition line as its span; its content is not reproduced here.\par\medskip
\ifdefined\hypertarget\hypertarget{sourcepointer-2}{}\fi
\noindent\textit{Source reference.} Original reference morphisms-lemma-identify-finite-type-points: sources/stacks-algebraic-stacks-morphisms/pinned-original-source.tex, source lines 3691--3691. Navigation pointer only: the counted text refers to the Stacks item morphisms-lemma-identify-finite-type-points, which belongs to another chapter of the project and is not part of this card. The pointer carries the target's own label as its number because this chapter's staged snapshot does not contain the other chapter's tag table; its content is not reproduced here.\par\medskip
\ifdefined\hypertarget\hypertarget{sourcepointer-3}{}\fi
\noindent\textit{Source reference.} Original reference spaces-morphisms-lemma-etale-locally-finite-presentation: sources/stacks-algebraic-stacks-morphisms/pinned-original-source.tex, source lines 3672--3672. Navigation pointer only: the counted text refers to the Stacks item spaces-morphisms-lemma-etale-locally-finite-presentation, which belongs to another chapter of the project and is not part of this card. The pointer carries the target's own label as its number because this chapter's staged snapshot does not contain the other chapter's tag table; its content is not reproduced here.\par\medskip
\ifdefined\hypertarget\hypertarget{sourcepointer-4}{}\fi
\noindent\textit{Source reference.} Original reference spaces-morphisms-lemma-finite-presentation-finite-type: sources/stacks-algebraic-stacks-morphisms/pinned-original-source.tex, source lines 305--305. Navigation pointer only: the counted text refers to the Stacks item spaces-morphisms-lemma-finite-presentation-finite-type, which belongs to another chapter of the project and is not part of this card. The pointer carries the target's own label as its number because this chapter's staged snapshot does not contain the other chapter's tag table; its content is not reproduced here.\par\medskip
\ifdefined\hypertarget\hypertarget{sourcepointer-5}{}\fi
\noindent\textit{Source reference.} Original reference spaces-morphisms-lemma-finite-type-points-morphism: sources/stacks-algebraic-stacks-morphisms/pinned-original-source.tex, source lines 3688--3688. Navigation pointer only: the counted text refers to the Stacks item spaces-morphisms-lemma-finite-type-points-morphism, which belongs to another chapter of the project and is not part of this card. The pointer carries the target's own label as its number because this chapter's staged snapshot does not contain the other chapter's tag table; its content is not reproduced here.\par\medskip
\ifdefined\hypertarget\hypertarget{sourcepointer-6}{}\fi
\noindent\textit{Source reference.} Original reference spaces-morphisms-lemma-identify-finite-type-points: sources/stacks-algebraic-stacks-morphisms/pinned-original-source.tex, source lines 3685--3685. Navigation pointer only: the counted text refers to the Stacks item spaces-morphisms-lemma-identify-finite-type-points, which belongs to another chapter of the project and is not part of this card. The pointer carries the target's own label as its number because this chapter's staged snapshot does not contain the other chapter's tag table; its content is not reproduced here.\par\medskip
\begin{lemma}
\label{lemma-point-finite-type}
Let $\mathcal{X}$ be an algebraic stack.
Let $x \in |\mathcal{X}|$. The following are equivalent:
\begin{enumerate}
\item There exists a morphism $\Spec(k) \to \mathcal{X}$
which is locally of finite type and represents $x$.
\item There exists a scheme $U$, a closed point $u \in U$, and a smooth
morphism $\varphi : U \to \mathcal{X}$ such that $\varphi(u) = x$.
\end{enumerate}
\end{lemma}

\begin{proof}
Let $u \in U$ and $U \to \mathcal{X}$ be as in (2). Then
$\Spec(\kappa(u)) \to U$ is of finite type, and $U \to \mathcal{X}$ is
representable and locally of finite type (by
Morphisms of Spaces,
Lemmas \ref{spaces-morphisms-lemma-etale-locally-finite-presentation} and
\ref{spaces-morphisms-lemma-finite-presentation-finite-type}).
Hence we see (1) holds by
Lemma \ref{lemma-composition-finite-type}.

\medskip\noindent
Conversely, assume $\Spec(k) \to \mathcal{X}$ is locally of finite type
and represents $x$. Let $U \to \mathcal{X}$ be a surjective smooth morphism
where $U$ is a scheme. By assumption
$U \times_\mathcal{X} \Spec(k) \to U$ is a morphism of algebraic
spaces which is locally of finite type. Pick a finite type point $v$ of
$U \times_\mathcal{X} \Spec(k)$ (there exists at least one, see
Morphisms of Spaces,
Lemma \ref{spaces-morphisms-lemma-identify-finite-type-points}).
By
Morphisms of Spaces,
Lemma \ref{spaces-morphisms-lemma-finite-type-points-morphism}
the image $u \in U$ of $v$ is a finite type point of $U$.
Hence by
Morphisms, Lemma \ref{morphisms-lemma-identify-finite-type-points}
after shrinking $U$ we may assume that $u$ is a closed point of $U$, i.e.,
(2) holds.
\end{proof}

\end{document}
